\section{Policy Translation and Ethical Use}

The screening framework should be used only as a preliminary diagnostic tool for supportive intervention. Its appropriate uses include identifying firms for outreach, offering diagnostic assessments, bundling finance and digital-capability support, and prioritizing follow-up research. It should not be used to deny credit, exclude firms from assistance, stigmatize businesses, or assign punitive risk labels.

Responsible deployment would require human review before any policy action. Firms should be able to understand why they were flagged, correct inaccurate information, and appeal decisions based on the score. The model should also be audited periodically across region, sector, firm size, ownership, and gender where data permit. These audits are necessary because a classifier that appears useful on average may still over-flag particular groups or under-detect need in others.

The governance principle is straightforward: fragility scores should expand access to support rather than restrict it. In practical terms, the score should trigger a conversation, not a final decision. It should be paired with qualitative assessment, field verification, and transparent eligibility rules. This is especially important in SME policy because risk scores can become self-fulfilling if they reduce access to finance or public support for firms already facing constraints.

Table~\ref{tab:policy_safeguards} below translates this principle into concrete policy-use boundaries. The table is intentionally conservative because the current score is constructed, only moderately predictive after leakage reduction, and not externally validated against observed firm distress.

\begin{table}[H]
\centering
\caption{Policy-Use Safeguards for SME Fragility Scores}
\label{tab:policy_safeguards}
\small
\resizebox{\textwidth}{!}{\begin{tabular}{llp{7.2cm}}
\toprule
Policy use case & Appropriate? & Required safeguard \\
\midrule
Diagnostic outreach & Yes & Use the score to invite voluntary assessment and support, with human review before follow-up action. \\
Credit denial & No & Prohibited use; the model is not validated for punitive or exclusionary financial decisions. \\
Grant prioritization & Conditional & Combine the score with field verification, transparent criteria, and an appeal mechanism. \\
Research sampling & Yes & Use only to identify cases for further study, with disclosure that the score is a constructed screening signal. \\
Automated eligibility & No & Not deployment-ready; eligibility decisions require independent validation, calibration, and governance review. \\
\bottomrule
\end{tabular}
}
\end{table}
